\documentclass[../../../include/open-logic-section]{subfiles}

\begin{document}

\olfileid[hi]{his}{set}{hilbertcurve}

\olsection{परिशिष्ट: हिल्बर्ट के अंतरिक्ष-पूर्ति वक्र}

अध्याय \olref[pathology]{sec} में हमने बताया था कि किसी रेखा और वर्ग पर
बिंदुओं की संख्या ठीक समान होने के कांटोर के प्रमाण
(\olref[his][set][cantorplane]{thm:cantorplane}) ने पियानो को यह पूछने के
लिए प्रेरित किया कि क्या कोई अंतरिक्ष-पूर्ति \emph{वक्र} हो सकता है। 1890
में उन्हें सकारात्मक उत्तर मिला। इस खंड में हम (कुछ अनौपचारिक ढंग से) समझाते
हैं कि हिल्बर्ट का अंतरिक्ष-पूर्ति वक्र (एक छोटे-से बदलाव सहित) कैसे बनाया
जाए।\footnote{अधिक कठोर व्याख्या के लिए \cite{Rose2010} देखें। बदलाव का अर्थ
नीचे के वक्रों के लाल भाग सम्मिलित करना है। इससे यह जाँचना कुछ सरल होता है कि
वक्र सतत है।}

हमें किसी फलन $h$ को फलनों $h_1$, $h_2$, $h_3$, \dots\@ के अनुक्रम की सीमा
के रूप में परिभाषित करना है। पहले हम निर्माण का वर्णन करेंगे। फिर दिखाएँगे कि
वह अंतरिक्ष की पूर्ति करता है। उसके बाद दिखाएँगे कि वह वक्र है।

हम $h$ का परिसर इकाई वर्ग $\unitsquare$ लेंगे। $h$ का हमारा पहला सन्निकटन,
अर्थात् $h_1$, यह है:
\begin{center}
	\begin{tikzpicture}
	\draw[gray] (0pt, 0pt) rectangle (80pt, 80pt);
	\draw[step = 40pt, gray, very thin] (0pt, 0pt) grid (80pt, 80pt);
	\draw[oldiagcolorC, thick] (20pt, 0)--(20pt, 20pt);
	\draw[oldiagcolorC, thick] (60pt, 0)--(60pt, 20pt);		
	\begin{scope}[xshift=20pt, yshift=20pt]
	\draw [black, thick, l-system={Hilbert curve, step=40pt, angle=90, axiom=L, order=1}]  lindenmayer system;
	\end{scope}
	\end{tikzpicture}
\end{center}
चीज़ों का हिसाब रखने के लिए हमने वर्ग पर $2 \times 2$ जाल रखा है। हम वक्र को
निचले बाएँ चतुर्थांश से आरंभ होकर ऊपरी बाएँ, फिर ऊपरी दाएँ और अंततः निचले
दाएँ जाते हुए सोच सकते हैं। निर्माण का दूसरा चरण, अर्थात् $h_2$, यह है:
\begin{center}
	\begin{tikzpicture}
	\draw[gray] (0pt, 0pt) rectangle (80pt, 80pt);
	\draw[step = 20pt, gray, very thin] (0pt, 0pt) grid (80pt, 80pt);
	\draw[oldiagcolorC, thick] (0pt, 10pt)--(10pt, 10pt);
	\draw[oldiagcolorC, thick] (70pt, 10pt)--(80pt, 10pt);
	\draw[thick, oldiagcolorE] (10pt, 30pt)--(10pt, 50pt); 
	\draw[thick, oldiagcolorE] (30pt, 50pt)--(50pt, 50pt); 
	\draw[thick, oldiagcolorE] (70pt, 50pt)--(70pt, 30pt); \begin{scope}[xshift=10pt, yshift=30pt]
	\draw [thick, rotate=270,black, l-system={Hilbert curve, step=20pt, angle=90, axiom=L, order=1}]  lindenmayer system;
	\end{scope}
	\begin{scope}[xshift=10pt, yshift=50pt]
	\draw [thick, black, l-system={Hilbert curve, step=20pt, angle=90, axiom=L, order=1}]  lindenmayer system;
	\end{scope}
	\begin{scope}[xshift=50pt, yshift=50pt]
	\draw [thick, black, l-system={Hilbert curve, step=20pt, angle=90, axiom=L, order=1}]  lindenmayer system;
	\end{scope}
	\begin{scope}[xshift=70pt, yshift=10pt]
	\draw [thick, rotate=90,black, l-system={Hilbert curve, step=20pt, angle=90, axiom=L, order=1}]  lindenmayer system;
	\end{scope}
	\end{tikzpicture}
\end{center}
भिन्न रंग यह समझाने में सहायता करेंगे कि $h_2$ कैसे बनाया गया। पहले हम
$h_1$ के गैर-!!{colorC} भाग की छोटे पैमाने की प्रतियाँ अपने वर्ग के निचले
बाएँ, ऊपरी बाएँ, ऊपरी दाएँ और निचले दाएँ भाग में रखते हैं (काले रंग में
खींची गईं)। फिर इन चार आकृतियों को (!!{colorE} रेखाओं से) जोड़ते हैं। अंत
में अपनी आकृति को वर्ग की सीमा से (!!{colorC} रेखाओं से) जोड़ते हैं।

अब $h_3$ लें। जैसे $h_2$, $h_1$ के गैर-लाल भाग की चार जुड़ी हुई छोटे पैमाने
की प्रतियों से बना था, वैसे ही $h_3$, $h_2$ के गैर-लाल भाग की चार छोटे
पैमाने की प्रतियों (काले रंग में खींची गईं) से बनता है, जिन्हें फिर
(!!{colorE} रेखाओं से) परस्पर और अंत में (!!{colorC} रेखाओं से) वर्ग की
सीमा से जोड़ा जाता है।
\begin{center}
	\begin{tikzpicture}
	\draw[gray] (0pt, 0pt) rectangle (80pt, 80pt);
	\draw[step = 10pt, gray, very thin] (0pt, 0pt) grid (80pt, 80pt);
	\draw[thick, oldiagcolorC](5pt, 0pt)--(5pt, 5pt); 
	\draw[thick, oldiagcolorC] (75pt, 0pt)--(75pt, 5pt); 
	\draw[thick, oldiagcolorE] (75pt, 45pt)--(75pt, 35pt);
	\draw[thick, oldiagcolorE] (5pt, 35pt)--(5pt, 45pt); 
	\draw[thick, oldiagcolorE] (35pt, 45pt)--(45pt, 45pt); 
	\draw[thick, oldiagcolorE] (75pt, 45pt)--(75pt, 35pt); \begin{scope}[xshift=5pt, yshift=35pt]
	\draw [rotate=270, thick, black, l-system={Hilbert curve, step=10pt, angle=90, axiom=L, order=2}]  lindenmayer system;
	\end{scope}
	\begin{scope}[xshift=5pt, yshift=45pt]
	\draw [thick, black, l-system={Hilbert curve, step=10pt, angle=90, axiom=L, order=2}]  lindenmayer system;
	\end{scope}
	\begin{scope}[xshift=45pt, yshift=45pt]
	\draw [thick, black, l-system={Hilbert curve, step=10pt, angle=90, axiom=L, order=2}]  lindenmayer system;
	\end{scope}
	\begin{scope}[xshift=75pt, yshift=5pt]
	\draw [thick, rotate=90,black, l-system={Hilbert curve, step=10pt, angle=90, axiom=L, order=2}]  lindenmayer system;
	\end{scope}
	\end{tikzpicture}
\end{center}
अब हमें $h_n$ से $h_{n+1}$ परिभाषित करने का सामान्य प्रतिरूप दिखता है। अंततः
हम इन क्रमागत फलनों $h_1$, $h_2$, \dots\@ की बिंदुवार सीमा लेकर \emph{स्वयं}
वक्र $h$ परिभाषित करते हैं। अर्थात् प्रत्येक $x \in \unitsquare$ के लिए:
\begin{align*}
	h(x) &= \lim_{n \rightarrow \infty} h_n(x)
\end{align*} 
अब हम दिखाते हैं कि यह वक्र अंतरिक्ष की पूर्ति करता है। जब हम वक्र $h_n$
खींचते हैं, तो $\unitsquare$ पर $2^n \times 2^n$ जाल रखते हैं। पाइथागोरस के
प्रमेय से प्रत्येक जाल-खाने के विकर्ण की लंबाई है:
\[
\sqrt{\left(\nicefrac{1}{2^{n}}\right)^2+\left(\nicefrac{1}{2^{n}}\right)^2} = 2^{(\frac{1}{2}-n)}
\]
और स्पष्टतः $h_n$ प्रत्येक जाल-खाने से होकर गुजरता है। अतः $\unitsquare$ का
प्रत्येक बिंदु $h_n$ के किसी बिंदु से \emph{अधिकतम}
$2^{(\frac{1}{2}-n)}$ दूरी पर है। अब $h$ को फलनों $h_1$, $h_2$, $h_3$,
\dots\@ की सीमा के रूप में परिभाषित किया गया है। इसलिए किसी भी बिंदु की $h$
से अधिकतम दूरी इससे दी जाती है:
\[
\lim_{n \rightarrow \infty} 2^{(\frac{1}{2}-n)} = 0.
\]
अर्थात् $\unitsquare$ का प्रत्येक बिंदु~$h$ से $0$ दूरी पर है। दूसरे शब्दों
में $\unitsquare$ का प्रत्येक बिंदु वक्र \emph{पर} स्थित है। अतः $h$
अंतरिक्ष की पूर्ति करता है!{}

यह दिखाना शेष है कि $h$ वास्तव में एक \emph{वक्र} है। इसे दिखाने के लिए हमें
संकल्पना परिभाषित करनी होगी। आधुनिक परिभाषा 1887 में जॉर्डन की दी हुई परिभाषा
पर आधारित है (अर्थात् पहला अंतरिक्ष-पूर्ति वक्र दिए जाने से केवल कुछ वर्ष
पहले की):

\begin{defn}
वक्र, $\unitline$ से $\Real^2$ में एक सतत प्रतिचित्रण है।
\end{defn}

यह काफ़ी अंतर्ज्ञान-सुलभ है: अंतर्ज्ञानतः वक्र एक ``सुचारु'' प्रतिचित्रण है
जो किसी विहित रेखा को समतल $\Real^2$ पर ले जाता है। हमारा फलन $h$ वास्तव में
$\unitline$ से $\Real^2$ में प्रतिचित्रण है। अतः हमें केवल दिखाना है कि $h$
सतत है। हमने \olref[limits]{sec} में $\epsilon$/$\delta$ संकेतन से सांतत्य
परिभाषित किया था। साधारण भाषा में हम यह स्थापित करना चाहते हैं:
\emph{यदि आप $\unitsquare$ में कोई बिंदु $p$ और परिशुद्धता का कोई इच्छित
स्तर $\epsilon$ निर्दिष्ट करें, तो हम $\unitline$ का ऐसा खुला खंड खोज सकते
हैं कि उस खुले खंड में कोई भी $x$ दिए जाने पर $h(x)$, $p$ से $\epsilon$ के
भीतर हो।}

अतः मान लें कि आपने $p$ और $\epsilon$ निर्दिष्ट किए हैं। इसका अर्थ वस्तुतः
$\unitsquare$ पर केंद्र $p$ और त्रिज्या $\epsilon$ वाला वृत्त खींचना है।
(वृत्त $\unitsquare$ के किनारे के बाहर निकल सकता है, किंतु इससे कोई अंतर
नहीं पड़ता।) अब याद करें कि फलन $h_n$ का वर्णन करते समय हमने $\unitsquare$
पर $2^n \times 2^n$ जाल खींचा था। स्पष्ट है कि $\epsilon$ कितना भी छोटा हो,
कोई ऐसा $n$ होता है कि $\unitsquare$ पर $2^n \times 2^n$ जाल का कोई एक खाना
केंद्र $p$ और त्रिज्या $\epsilon$ वाले वृत्त के पूर्णतः भीतर स्थित हो।

अतः वह $n$ लें और $I$, $\unitline$ का वह सबसे बड़ा खुला भाग हो जिसे $h_n$
पूरी तरह संबद्ध जाल-खाने में प्रतिचित्रित करता है। (स्पष्ट है कि $(a,b)$
अस्तित्व में है, क्योंकि हम पहले ही बता चुके हैं कि $h_n$, $2^n\times 2^n$
जाल के प्रत्येक खाने से होकर गुजरता है।) अब इतना दिखाना पर्याप्त है कि जब
भी $x \in I$, बिंदु $h(x)$ उसी जाल-खाने में स्थित होता है। और \emph{यह}
दिखाने के लिए इतना दिखाना पर्याप्त है कि किसी भी $m > n$ के लिए $h_m(x)$
उसी जाल-खाने में है। किंतु यह स्पष्ट है। यदि हम देखें कि $m > n$ के लिए
$h_m$ के साथ क्या होता है, तो पाते हैं कि इकाई अंतराल का ठीक ``वही भाग'' उसी
जाल-खाने में प्रतिचित्रित होता है; बस हम उसे उस क्षेत्र में उत्तरोत्तर अधिक
खींचे और लहरदार ढंग से प्रतिचित्रित करते हैं।

\end{document}
